\section*{第 7 讲\quad 链复形与 Herbrand 商}

上一讲我们对循环群定义了 Tate 上同调群 $\HH^0(G,M)$ 与 $\HH^1(G,M)$。回忆：若 $G=\Z/n\Z$、生成元为 $\sigma$，则一个 $G$-模就是带自同构 $\sigma:M\xrightarrow{\sim}M$ 且满足 $\sigma^n=\mathrm{id}_M$ 的交换群 $M$。我们的主要例子来自域扩张 $L/K$，其中 $\Gal(L/K)=G$，此时 $L$ 与 $L^\times$ 都是 $G$-模。于是
\[
\HH^0(G,M):=M^G/N(M)=\ker(1-\sigma)\Big/\Big\{\sum_{i=0}^{n-1}\sigma^i m:m\in M\Big\},
\]
\[
\HH^1(G,M):=\ker(N)/(1-\sigma),
\]
因为 $\ker(1-\sigma)$ 中的元素在 $\sigma$ 作用下不动，从而在 $G$ 作用下不动。我们的目标是在上述例子中证明 $\#\HH^0=n$，所用的工具是长正合列。

我们看到，若
\[
0\to M\to E\to N\to 0
\]
是 $G$-模的短正合列（即 $M,E,N$ 都是带 $n$ 阶自同构的交换群，这些自同构与上述映射相容，且 $N=E/M$，于是 $M$ 在 $N$ 的自同构下不动），则得到长正合列
\[
\HH^0(G,M)\to\HH^0(G,E)\to\HH^0(G,N)\xrightarrow{\delta}\HH^1(G,M)\to\HH^1(G,E)\to\HH^1(G,N),
\]
其中边缘映射 $\delta$ 把 $x\in\HH^0(N)=N^G/N(N)$ 提升为 $\tilde x\in E$，使得 $(1-\sigma)\tilde x\in\ker(N)\subseteq M$，从而给出 $\HH^1(G,M)$ 中的一个类。

现在定义第二个边缘映射
\begin{equation}\tag{7.1}
\HH^1(G,M)\to\HH^1(G,E)\to\HH^1(G,N)\xrightarrow{\partial}\HH^0(G,M)\to\HH^0(G,E)\to\HH^0(G,N),
\end{equation}
它把 $x\in\HH^1(G,N)$ 提升为元素 $\tilde x\in E$。此时 $N(\tilde x)=\sum_{i=0}^{n-1}\sigma^i\tilde x\in M^G$，因为它被 $1-\sigma$ 杀死，于是它定义了 $\HH^0(G,M)$ 中的一个类。我们验证：

\textbf{断言 7.1.} 边缘映射 $\partial$ 是良定义的。

\textit{证明.} 若 $\tilde x'$ 是 $x$ 的另一个提升，则 $\tilde x-\tilde x'\in M$（因为 $N=E/M$），从而 $\sum_{i=0}^{n-1}\sigma^i(\tilde x-\tilde x')\in N(M)$，在 $\HH^0(G,M)$ 中被消掉。\hfill$\square$

\textbf{断言 7.2.} (7.1) 中的序列正合。

\textit{证明.} 若 $x\in\HH^1(G,E)$，则 $N(x)=0$，故 $\partial(x)=0$ 于 $\HH^0(G,M)$。若 $x\in\ker(\partial)$，则对 $x$ 的某个提升 $\tilde x\in E$ 有 $N(\tilde x)=0$，而 $x$ 正是 $\tilde x$ 的像。

若 $x\in\HH^1(G,N)$、取提升 $\tilde x\in E$，则由定义 $\partial(x)=N(\tilde x)$ 在 $\HH^0(G,E)$ 中为零。若 $x\in\HH^0(G,M)$ 在 $\HH^0(G,E)$ 中为零，则 $x\in N(E)$，从而 $x\in\im(\partial)$。\hfill$\square$

于是我们得到循环群 Tate 上同调的「$2$-周期」正合列，这引出如下定义。

\textbf{定义 7.3.} 对每个 $i\in\Z$（正、负均可），定义
\[
\HH^i(G,M):=
\begin{cases}
\HH^0(G,M), & i\equiv 0 \pmod 2,\\
\HH^1(G,M), & i\equiv 1 \pmod 2.
\end{cases}
\]
这一良好性质对非循环群不成立，因此我们常设法把上同调化归到循环群的情形。

作为另一种写法，置
\begin{equation}\tag{7.2}
\cdots\xrightarrow{\sum_{i=0}^{n-1}\sigma^i}M\xrightarrow{1-\sigma}M\xrightarrow{\sum_{i=0}^{n-1}\sigma^i}M\xrightarrow{1-\sigma}\cdots,
\end{equation}
并注意它构成我们将要称为链复形的东西：

\textbf{定义 7.4.} 链复形 $X_\bullet$ 是一列
\[
\cdots\xrightarrow{d_{-2}}X_{-1}\xrightarrow{d_{-1}}X_0\xrightarrow{d_0}X_1\xrightarrow{d_1}X_2\xrightarrow{d_2}\cdots,
\]
% [原文疑误：原文作 $\ker(d_{i+1})\supset\ker(d_i)$，按 $d_{i+1}d_i=0$ 应为 $\mathrm{im}(d_i)\subseteq\ker(d_{i+1})$；译文径改]
满足对一切 $i\in\Z$ 有 $d_{i+1}d_i=0$（也就是说 $\im(d_i)\subseteq\ker(d_{i+1})$，但不必像正合列那样取等号）。于是定义 $X_\bullet$ 的第 $i$ 个上同调为
\[
H^i(X_\bullet):=\ker(d_i)/\im(d_{i-1}).
\]
因此，长正合列是一种特殊的链复形。注意 (7.2) 满足该定义，因为
\[
(1-\sigma)\sum_{i=0}^{n-1}\sigma^i x=\sum_{i=0}^{n-1}\sigma^i x-\sum_{i=0}^{n-1}\sigma^{i+1}x=Nx-Nx=0,
\]
且两个映射显然交换。Tate 上同调群于是就是该链复形的上同调，由此可以清楚地看出它们具有 $2$-周期性。

\textbf{定义 7.5.} $G$-模 $M$ 的 \textbf{Herbrand 商}或\textbf{欧拉示性数}为
\[
\chi(M):=\frac{\#\HH^0(G,M)}{\#\HH^1(G,M)},
\]
它只在分子与分母都有限时才有定义。

此定义推广了我们先前对平凡 $G$-模的讨论：此时 $\HH^0(G,M)=M/n$、$\HH^1(G,M)=M[n]$；不过要注意，从偶上同调到奇上同调的边缘映射为零。

\textbf{引理 7.6.} 设
\[
0\to M\to E\to N\to 0
\]
是 $G$-模的短正合列。若 $\chi$ 对三个 $G$-模中的两个有定义，则它对三个都有定义，此时
\[
\chi(M)\cdot\chi(N)=\chi(E).
\]

\textit{证明.} 作出长正合列
\[
0\to\ker(\alpha)\to\HH^0(M)\xrightarrow{\alpha}\HH^0(E)\to\HH^0(N)\xrightarrow{\delta}\HH^1(M)\to\HH^1(E)\xrightarrow{\beta}\HH^1(N)\to\coker(\beta)\to 0 .
\]
由于第二个边缘映射给出正合列
\[
\HH^1(E)\xrightarrow{\beta}\HH^1(N)\xrightarrow{\partial}\HH^0(M)\xrightarrow{\alpha}\HH^0(E),
\]
我们有
\[
\ker(\alpha)=\im(\partial)=\HH^1(N)/\ker(\partial)=\HH^1(N)/\im(\beta)=\coker(\beta).
\]
应用引理 6.4 并消去 $\#\ker(\alpha)$ 与 $\#\coker(\beta)$，即得欲证之结论（与引理 6.7 的情形相同）。\hfill$\square$

关于有限性，插入一段简短的讨论：

\textbf{断言 7.7.} 群 $\HH^0(G,M)$ 与 $\HH^1(G,M)$ 都是 $n$-挠的。

\textit{证明.} 设 $x\in M^G$。则 $N(x)=\sum_{i=0}^{n-1}\sigma^i x=\sum_{i=0}^{n-1}x=nx$。于是 $nx\in N(M)$，故 $\HH^0(G,M)$ 是 $n$-挠的。现设 $x\in\ker(N)$。则
\[
nx=nx-Nx=\sum_{i=1}^{n}(1-\sigma^i)x=(1-\sigma)\sum_{i=1}^{n}(1+\cdots+\sigma^{i-1})x,
\]
因此 $nx\in(1-\sigma)M$，$\HH^1(G,M)$ 同样是 $n$-挠的。\hfill$\square$

于是，$\HH^0(G,M)$ 与 $\HH^1(G,M)$ 的有限生成性蕴含有限性。现在回忆我们的目标：对局部域的循环 $n$ 次扩张 $L/K$，证明 $\#\HH^0(L^\times)=n$。我们有如下更精细的论断：

\textbf{断言 7.8.} 保持上述设定，则
\begin{enumerate}
\item[(1)] $\HH^1(L^\times)=0$（从而 $\chi(L^\times)=\#\HH^0(L^\times)$）；
\item[(2)] $\chi(\OK_L^\times)=1$；
\item[(3)] $\chi(L^\times)=n$。
\end{enumerate}

\textit{证明.} 先证 (2) 蕴含 (3)。有正合列
\[
1\to\OK_L^\times\to L^\times\xrightarrow{v}\Z\to 0,
\]
其中 $v$ 表示赋值。于是由引理 7.6，
\[
\chi(L^\times)=\chi(\OK_L^\times)\cdot\chi(\Z)=1\cdot n=n
\]
（用到 (2)），这里注意
\[
\HH^0(\Z)=\Z^G/N\Z=\Z/n\Z,\qquad
\HH^1(\Z)=\ker(N)/(1-\sigma)=0 .
\]
下面证明 (2)。

\textbf{引理 7.9.} 若 $M$ 是有限 $G$-模，则 $\chi(M)=1$。

\textit{证明.} 有正合列
\[
0\to M^G\xrightarrow{1-\sigma}\ker(N)\to\HH^1(G,M)\to 0,
\]
\[
0\to\ker(N)\to M\xrightarrow{\sum_{i=0}^{n-1}\sigma^i}M^G\to\HH^0(G,M)\to 0,
\]
于是由引理 7.6，
\[
\#\ker(N)\cdot\#M^G=\#M\cdot\#\HH^0(G,M),\qquad
\#M^G\cdot\#\ker(N)=\#M\cdot\#\HH^1(G,M),
\]
故 $\#\HH^0(G,M)=\#\HH^1(G,M)$，即 $\chi(M)=1$，如所欲证。\hfill$\square$

类似地有 $\chi(\OK_L)=1$，其中把 $\OK_L$ 视为加法群。事实上，更容易建立的说法是 $\chi(L)=1$。直观上，这是因为我们工作在 $p$ 进数上，一切都必须是 $\Q$-向量空间，故 $n$ 可逆；而由断言 7.7，我们的上同调群都是 $n$-挠的，因此两个上同调群都必须消没，从而 $\chi(L)=1$。

由正规基定理，若 $L/K$ 是有限 Galois 扩张，则作为 $K[G]$-模有
\[
L\simeq\prod_{g\in G}K=K[G],
\]
其中 $G$ 通过置换坐标作用。这是因为 $K$ 的作用（数乘；$L$ 是 $K$-向量空间）与 $G$ 的作用（$G$ 作为 $K$-向量空间作用于 $L$）交换，故 $L$ 上有 $K[G]$-作用。

\textbf{断言 7.10.} 设 $A$ 是任意交换群，$A[G]:=\prod_{g\in G}A$，$G$ 通过置换坐标作用于此。若 $G$ 循环，则
\[
\HH^0(G,A[G])=\HH^1(G,A[G])=0 .
\]

\textit{证明.} 我们把断言改述如下：设 $R$ 是交换环，则 $R[G]$ 通过乘法运算
\[
\Big(\sum_{g\in G}x_g[g]\Big)\Big(\sum_{h\in G}y_h[h]\Big):=\sum_{g,h\in G}x_gy_h[gh]
\]
成为（可能非交换的）$R$-代数，其中我们用 $[h]\in\prod_{g\in G}R$ 记第 $h$ 个坐标为 $1$、其余坐标为 $0$ 的元素。于是 $R[G]$-模等价于带同态 $G\to\operatorname{Aut}_R(M)$ 的 $R$-模。特别地，$\Z[G]$-模等价于 $G$-模。

现在 $\HH^0(G,A[G])=A[G]^G/N$，其中 $A[G]^G$ 等同于对角嵌入的 $A\subset\prod_{g\in G}A$，而 $N((a,0,\dots,0))=\sum_{g\in G}a[g]$ 等于 $A$ 的对角嵌入，故 $\HH^0(G,A[G])=0$。

同理，$\HH^1(G,A[G])=\ker(N)/(1-\sigma)$，且
\[
A[G]\supseteq\ker(N)=\Big\{\sum_{g\in G}a_g[g]\in A[G]:\sum_{g\in G}a_g=0\Big\}.
\]
现在，一般元素可写成 $\sum_{i=0}^{n-1}a_i[\sigma^i]$，并对每个 $i$ 选取 $b_i$ 使得
\[
(1-\sigma^{n-i})a_i=(1-\sigma)b_i .
\]
于是
\[
(1-\sigma)\sum_{i=0}^{n-1}b_i[\sigma^i]
=\sum_{i=0}^{n-1}(1-\sigma^{n-i})a_i[\sigma^i]
=\sum_{i=0}^{n-1}a_i[\sigma^i]-\sum_{i=0}^{n-1}a_i[1]
=\sum_{i=0}^{n-1}a_i[\sigma^i],
\]
故 $\ker(N)\subset(1-\sigma)A[G]$，从而 $\HH^1(G,A[G])=0$，如所欲证。\hfill$\square$

由此可见，用这种方式得不到有趣的 Tate 上同调。

现在我们回到 $\chi(\OK_L)=1$ 的证明。困难在于正规基定理对 $\OK_L$ 并不适用：虽然 $L=K[G]$，我们却不一定有 $\OK_L\simeq\OK_K[G]$。

不过，$\OK_L$ 确实存在具有正规基的开子群。取 $L/K$ 的一组正规基 $\{e_1,\dots,e_n\}$。对充分大的 $N$，有 $\pi^Ne_1,\dots,\pi^Ne_n\in\OK_L$，其中 $\pi$ 是 $L$ 的素元；于是它们自由张成 $\OK_L$ 的某个开子群。由于这个子群（记作 $\Gamma$）具有有限指数，由 (6.2) 得
\[
\chi(\OK_L)=\chi(\Gamma)=\chi(\OK_K[G])=1 .
\]

为证 $\chi(\OK_L^\times)=1$（更完整的证明将在下一讲给出），注意 $\OK_L^\times\supseteq\Gamma\simeq\OK_L^+$ 是 $G$-等价的，其中 $\Gamma$ 是开子群（这一事实的证明用到 $p$ 进指数映射）。于是 $\chi(\OK_L^\times)=\chi(\Gamma)=1$，如所欲证。\hfill$\square$

\textbf{注 7.11.} 本课程中，所有环与模均假定含幺。
